% What this holds: Section Demographic Representation and its lead-in paragraph.
% Read by arlington-bsap.tex. Prose lives here; formatting lives in the wrapper and style/.
\section{Demographic Representation}\label{sec:who}

In September 2019 the Arlington County Board adopted an equity resolution
committing the County to collect and disaggregate data on disparities and to ask
of its decisions who benefits, who is burdened and who is
missing.\autocite{arlingtoncb2019equity} This report turns that last question on
the Board itself. It reads the Board's history through the resolution, asking who
has held the Board's seats and who has not: how many seats there were, who sat in
them, and how that has changed for women, for residents of color and for the
young.

Of the \membersTotal{} people who have served on the Board since 1870,
\whiteMenMembers{} have been white men. Table~\ref{tab:subsets} lists the
other \notWhiteMenMembers{}, in the order they first took a seat: 13 women
and 12 members of color, Tannia Talento among both.

% waits on: member-demographics-lists
\begin{table}[b]
  \refstepcounter{table}\label{tab:subsets}
  {\raggedright\MakeUppercase{Table \thetable}\par}
  {\raggedright\bfseries Women and Members of Color on the Board\par}
  \medskip
  \begingroup\footnotesize\setlength{\tabcolsep}{8pt}
  \begin{tabular}{Q}
  \toprule
  \rosterhead
  \midrule
  \rosterinput members_roster_subsets.tex
  \bottomrule
  \end{tabular}\endgroup
  \par\smallskip
  \figurenotes{
    Notes: \rostermarks{} This table is a subset of the full roster, Table~\ref{tab:roster} in the Data Appendix.}{}
\end{table}

\clearpage
